\versiondate{3.3.03}
%\def\vtmp#1{#1}
%\def\vtmpb#1{(#1)}
%\def\vtmp#1{}
\def\vtmpb#1{}
\def\newvolume#1{}
     
\newif\iflargelogo
     
\def\pagereference#1#2{\hfill\hbox{\ifresultsonly#2\qquad \hskip 1truein
                                                 \else#1\qquad \fi}}
\long\def\section#1#2#3#4#5#6{\ifresultsonly
    \hskip1truein\qquad #1 #2\pagereference{#4}{#5}
  \else\vskip 2pt plus 1pt minus 0pt
    \vbox{\qquad #1 #2\pagereference{#4}{#5}\par
    \abstr{#6}
    }\vskip -1pt plus 1pt minus 0pt
  \fi}
\long\def\chapintrosection#1#2#3{\ifresultsonly
    \hskip1truein\qquad Pendahuluan\pagereference{#2}{#3}
  \else\vskip 1pt plus 1pt minus 0pt
    \qquad Pendahuluan\pagereference{#2}{#3}
    \vskip 0pt plus 1pt minus 0pt
  \fi}
     
\def\abstr#1{\ifresultsonly{}
   \else{\advance\leftskip by\parindent
     \advance\rightskip by\parindent
     \advance\rightskip by\parindent
     {\advance\leftskip by\parindent
      \advance\rightskip by\parindent
     {\advance\leftskip by\parindent
     {\advance\leftskip by\parindent
     {\vskip1pt plus 0pt minus0pt
     \parindent=0pt
     \smallerfonts #1
     \vskip1pt plus 0pt minus0pt }}}}}\fi}
     
\Loadfourteens
\hyphenation{Dede-kind}
     
\gdef\topparagraph{}
\gdef\bottomparagraph{}
     
\vbox{\vskip 2truein
     
\centerline{\fourteenbf TEORI UKURAN}
}%end of vbox
     
\bigskip
     
\centerline{\fourteenbf Jilid 2}
     
\ifresultsonly
   \bigskip
     
   \centerline{\fourteenit (hanya hasil)}
   \fi
     
\vskip 2truein
     
\centerline{\fourteenrm D.H.Fremlin}
     
\vskip 2truein
     
\largelogofalse
\input mtlogo
     
\vfill\eject
     
\vbox{\vskip 2truein
     
\noindent Karya lain oleh penulis yang sama:
     
{\it Topological Riesz Spaces and Measure Theory}, Cambridge University
Press, 1974.
     
{\it Consequences of Martin's Axiom}, Cambridge University Press, 1982.
}%end of vbox
     
\bigskip
     
\noindent Jilid pendamping untuk jilid ini:
     
{\it Measure Theory}, vol.\ 1, Torres Fremlin, 2000;
     
{\it Measure Theory}, vol.\ 3, Torres Fremlin, 2002;
     
{\it Measure Theory}, vol.\ 4, Torres Fremlin, 2003;

{\it Measure Theory}, vol.\ 5, Torres Fremlin, 2008.
     
\vfill
     
\Centerline{\it Edisi pertama Mei 2001}
     
\Centerline{\it Edisi sampul keras Januari 2010}
     
\Centerline{\it Cetakan kedua April 2016}
     
\vfill\eject
\def\Loadtwenties{
      \font\twentyrm=cmr10 scaled \magstep4
      \font\twentybf=cmbx10 scaled \magstep4
      \font\twentyit=cmsl10 scaled \magstep4}
     
\Loadtwenties
     
\vbox{\vskip 1.5truein
     
\centerline{\twentybf TEORI UKURAN}
}%end of vbox
     
\bigskip\bigskip
     
\centerline{\twentybf Jilid 2}
     
\bigskip\bigskip
     
\centerline{\twentyrm Landasan yang Luas}
     
\ifresultsonly
   \bigskip\bigskip
     
   \centerline{\twentyit (hanya hasil)}
   \fi
     
\vskip 1truein
\vfill
     
\centerline{\twentyrm D.H.Fremlin}
     
\bigskip\bigskip
     
\centerline{\fourteenit Profesor Emeritus Matematika,
University of Essex}
     
\vskip 1truein
\vfill
     
\largelogotrue
\input mtlogo
     
\eject
     
\vbox{\vskip 5truecm
     
\Centerline{\fourteenit Dipersembahkan oleh Penulis}
     
\Centerline{\fourteenit kepada Penerbit}
}%end of vbox
     
\vfill
     
\Centerline{Buku ini dapat dipesan dari percetakan,
{\tt http://www.lulu.com/buy}.}

\vfill
     
\bigskip
     
\noindent Pertama kali diterbitkan pada tahun 2001
     
oleh Torres Fremlin, 25 Ireton Road, Colchester CO3 3AT, England
     
\medskip
     
\noindent\copyright\ D.H.Fremlin 2001
     
\noindent Hak D.H.Fremlin untuk diakui sebagai penulis karya ini telah
ditegaskan sesuai dengan Copyright, Designs and Patents Act 1988.
Karya ini diterbitkan berdasarkan ketentuan Design Science License sebagaimana
dipublikasikan di {\tt http://www.gnu.org/licenses/dsl.html}.   Untuk berkas
sumber, lihat {\tt http://www.essex.\discretionary{}
{}{}ac.uk/maths/people/\discretionary{}
{}{}fremlin\discretionary{}{}{}/mt2.2016/index.htm}.
     
%No commercial use may be made of any part of this work without
%permission in writing from the publisher.
     
\medskip
     
\noindent Klasifikasi Library of Congress QA312.F72
     
\medskip
     
\noindent Klasifikasi AMS 2010 28-01
     
\medskip
     
\noindent ISBN 978-0-9538129-7-4
     
\medskip
     
\noindent Ditata huruf dengan \AmSTeX
     
\medskip
     
\noindent Dicetak oleh Lulu.com
     
\eject
     
\pageno=5
     
\vbox{\ifresultsonly\vskip1truein\fi
     
\centerline{\bf Daftar Isi}}
     
\bigskip
     
Pendahuluan Umum \pagereference{9}{}
     
\medskip
     
\input mt02
     
\vfill\eject
     
\gdef\topparagraph{}
\gdef\bottomparagraph{{\it Pendahuluan umum}}
     
\def\volumename{Landasan yang luas}
     
\bigskip
     
\noindent{\bf Pendahuluan Umum} Dalam risalah ini saya bermaksud memberikan
uraian menyeluruh tentang teori ukuran abstrak modern, beserta beberapa
gambaran mengenai penerapan utamanya.   Dua jilid pertama disusun pada
tingkat pengantar;  keduanya ditujukan bagi mahasiswa yang memiliki
landasan kokoh dalam konsep-konsep analisis riil, tetapi mungkin memiliki
pengetahuan terperinci yang agak terbatas.    Seiring berjalannya buku ini,
tingkat kecanggihan dan keahlian yang dituntut akan meningkat;  dengan
demikian, untuk jilid tentang ruang ukur topologis, penguasaan topologi umum
akan dianggap sudah ada.
Penekanan di seluruh karya ini terletak pada gagasan-gagasan matematika yang
terlibat, yang dalam bidang ini sebagian besar terdapat dalam rincian
pembuktian.
     
Niat saya adalah agar buku ini dapat digunakan baik sebagai pengantar
awal tentang bidang ini maupun sebagai karya acuan.   Demi tujuan pertama,
saya berusaha membatasi gagasan dalam jilid-jilid awal pada hal-hal yang
benar-benar esensial bagi pengembangan teorema-teorema dasar.
Demi tujuan kedua, saya berusaha mengungkapkan gagasan-gagasan ini dalam
keumuman alaminya secara utuh, dan khususnya saya berhati-hati untuk tidak
menyiratkan pembatasan yang tidak perlu pada cakupan penerapannya.   Tentu
saja, sampai taraf tertentu prinsip-prinsip ini saling bertentangan.
Meskipun demikian, saya mendapati bahwa hampir sepanjang waktu keduanya
hampir dapat diselaraskan, {\it asalkan} saya membiarkan adanya sejumlah
pengulangan.   Misalnya, tepat pada awal pembahasan muncul pertanyaan:
haruskah ukuran Lebesgue dikembangkan terlebih dahulu pada garis bilangan
riil, lalu pada ruang berdimensi lebih tinggi, ataukah sebaiknya langsung
membahas kasus multidimensi?   Saya percaya tidak ada satu jawaban yang benar
untuk pertanyaan ini.   Kebanyakan mahasiswa akan menganggap kasus satu
dimensi lebih mudah, sehingga kasus itu tampak lebih sesuai untuk pengantar
pertama, sebab bahkan dalam kasus tersebut masalah teknisnya dapat terasa
berat.   Namun, setiap mahasiswa teori ukuran tentu harus pada tahap yang
cukup awal
memahami luas dan volume Lebesgue serta panjang;
dan dengan perumusan yang tepat, kasus multidimensi berbeda dari kasus satu
dimensi hanya dalam satu definisi dan satu lema (yang cukup berbobot).
Karena itu, yang saya lakukan adalah menuliskan keduanya (\S\S114-115).
Dengan semangat yang sama, saat menyusun latihan saya tidak menahan diri
hanya karena banyak hasil yang saya minta mahasiswa cari akan muncul
dalam bab-bab berikutnya;  saya percaya bahwa di seluruh matematika, seseorang
lebih berpeluang memahami suatu teorema jika sebelumnya ia telah mencoba
sendiri sesuatu yang serupa.
     
Rencana karya ini adalah sebagai berikut:
     
\medskip
     
\qquad\qquad Jilid 1:  Minimum yang Tak Tereduksi
     
\qquad\qquad Jilid 2:  Landasan yang Luas
     
\qquad\qquad Jilid 3:  Aljabar Ukuran
     
\qquad\qquad Jilid 4:  Ruang Ukur Topologis
     
\qquad\qquad Jilid 5:  Teori Ukuran dalam Kerangka Teori Himpunan.
     
\medskip
     
\noindent Jilid 1 ditujukan bagi mereka yang belum memiliki pengetahuan
sebelumnya tentang teori ukuran, tetapi menguasai teknik-teknik elementer
analisis riil.  Saya berharap jilid ini berguna bagi mahasiswa program sarjana
yang untuk pertama kalinya mempelajari ukuran Lebesgue.    Jilid 2 bertujuan
memaparkan beberapa hasil fundamental teori ukuran murni (teorema
Radon-Nikod\'ym, teorema Fubini), tetapi juga memberikan pengantar singkat
ke beberapa penerapan teori ukuran yang terpenting (teori probabilitas,
analisis Fourier).   Walaupun saya ingin percaya bahwa sebagian besar isinya
ditulis pada tingkat yang dapat dijangkau oleh siapa pun yang telah menguasai
isi Jilid 1, saya sendiri tidak akan cukup berani untuk mencoba membahas
seluruhnya dalam mata kuliah tingkat sarjana, meskipun saya tentu akan
berusaha memasukkan beberapa bagiannya.   Jilid 3 dan 4 disusun pada tingkat
yang agak lebih tinggi, sesuai untuk mata kuliah pascasarjana;  sedangkan
Jilid 5 mengandaikan penguasaan luas atas banyak bagian analisis dan
teori himpunan.
     
Ada sebuah catatan penafian yang patut saya sampaikan di tempat yang dapat
Anda lihat sebelum terlanjur membayar buku ini.   Saya tidak berusaha
memaparkan sejarah bidang ini.   Ini bukan karena saya menganggap sejarahnya
tidak menarik atau tidak penting;  alasannya justru karena saya tidak yakin
dapat mengatakan sesuatu yang tidak akan sangat menyesatkan.
Bahkan, saya sangat meragukan segala sesuatu yang pernah saya baca mengenai
sejarah gagasan.   Jadi, sementara saya dengan senang hati
menghormati nama Lebesgue, Kolmogorov, dan Maharam di tempat-tempat yang
kurang lebih sesuai, dan saya berusaha memasukkan ke dalam bibliografi
karya-karya yang saya sendiri gunakan, pembahasan rinci saya serahkan kepada
mereka yang lebih berani dan lebih memenuhi syarat daripada saya.
     
Untuk sementara,
setidaknya, pencetakan akan dilakukan dalam tiras kecil.   Saya berharap para
pembaca akan giat menyampaikan komentar mengenai kesalahan dan kekurangan,
karena keduanya seharusnya dapat diperbaiki dengan relatif cepat.   Salah satu
konsekuensi yang tak terelakkan ialah bahwa rujukan paragraf dapat lekas
kedaluwarsa.   Saya akan sangat tersanjung jika ada yang memilih mengandalkan
buku ini sebagai sumber materi dasar;  dan saya bersedia mencoba
mempertahankan konkordansi untuk rujukan semacam itu, dengan menunjukkan di
mana hasil-hasil yang berpindah untuk sementara berlabuh, jika para penulis
memberi saya salinan makalah yang menggunakannya.   Dalam konkordansi untuk 
jilid ini, Anda akan menemukan catatan mengenai butir-butir yang telah dirujuk
dalam jilid lain karya ini yang sudah diterbitkan.
     
Saya menyebutkan beberapa hal kecil mengenai tata letak materi.
Kebanyakan bagian ditutup dengan daftar `latihan dasar' dan
`latihan lanjutan', yang saya harap pada umumnya bersifat mendidik dan
sesekali menghibur.   Anda dan pengajar Anda, jika ada, harus memutuskan
berapa banyak yang perlu Anda coba, karena tidak ada dua mahasiswa yang
memiliki kebutuhan persis sama.   Saya menandai latihan yang menurut saya
sangat penting dengan $\pmb{>}$.   Namun, meskipun Anda mungkin tidak perlu
menuliskan penyelesaian untuk semua `latihan dasar', jika Anda meragukan
kemampuan Anda untuk melakukannya, anggaplah ini sebagai peringatan untuk
sedikit memperlambat langkah.   `Latihan lanjutan' memiliki tingkat kesulitan
yang tak berbatas,
dan satu-satunya kesamaan di antaranya adalah anggapan bahwa masing-masing
memiliki sekurang-kurangnya satu penyelesaian berdasarkan gagasan yang sudah
diperkenalkan.   Sesekali saya menambahkan sebuah `masalah' terakhir,
suatu pertanyaan yang jawabannya tidak saya ketahui dan yang tampaknya muncul
secara alami dalam perkembangan karya ini.
     
Dorongan untuk menulis risalah ini sebagian besar berasal dari keinginan
untuk menyajikan uraian terpadu mengenai bidang ini.   Karena itu, rujuk silang
berlimpah dan menjangkau banyak bagian.  Agar dapat merujuk dengan bebas ke
seluruh
teks, saya telah memilih sistem rujukan yang memberikan kode yang sama kepada
sebuah paragraf dari mana pun paragraf itu dirujuk.   Dengan demikian, 132E
adalah paragraf kelima dalam bagian kedua bab ketiga Jilid 1, dan dirujuk
dengan nama itu di seluruh karya.
Perlu saya tekankan bahwa rujuk silang dimaksudkan untuk membantu pembaca,
bukan mengalihkan perhatiannya.   Jangan menganggap sisipan `(121A)' sebagai
perintah, atau bahkan anjuran, untuk mengambil Jilid 1 dari rak lalu mencari
\S121.   Jika Anda merasa puas dengan argumen sebagaimana adanya,
terlepas dari rujukannya, lanjutkanlah.   Namun, jika saya tampak
melakukan lompatan yang agak besar, atau notasinya tiba-tiba menjadi
tidak jelas, rujuk silang setempat dapat membantu Anda mengisi celahnya.
     
Setiap jilid akan memiliki lampiran `fakta-fakta berguna', tempat saya
memaparkan materi yang digunakan di suatu bagian jilid tersebut dan yang
menurut saya tidak dapat dianggap sudah diketahui.   Biasanya penataan materi
dalam lampiran ini diarahkan secara sangat khusus pada penerapan tertentu
yang saya maksudkan, dan agaknya tidak dapat menjadi pengganti yang memadai
bagi pembahasan konvensional mengenai topik-topik yang disinggung.
Selain itu, gagasan-gagasan tersebut mungkin hanya diperlukan pada kesempatan
yang jarang dan terpisah-pisah.   Karena itu, sebagai aturan umum saya
menyarankan Anda mengabaikan lampiran sampai Anda mempunyai alasan langsung
untuk menduga bahwa suatu bagiannya mungkin berguna bagi Anda.
     
Selama proses panjang pematangan proyek ini saya telah dibantu oleh banyak
orang, dan saya berharap teman serta kolega saya akan senang ketika mereka
mengenali gagasan mereka yang tersebar di halaman-halaman berikut.   Namun,
saya terutama berterima kasih kepada mereka yang telah bersusah payah membaca
draf-draf terdahulu dan mengomentari bagian yang tidak jelas serta kesalahan.
     
\wheader{{\it Pendahuluan Jilid 2}}{18}{6}{6}{48pt}
     
\input mt2
     
